\chapter{Допуск determinant-line свёртки горизонтальной фазы}
\label{ch:v8-horizontal-phase-determinant-line-admission-gate}

Для инцидентности $A_0:\mathbb C^{11}\to\mathbb C^{10}$ определим
кофакторный вектор максимальных миноров
\begin{equation}
 c_j(A)=(-1)^j\det A_{\widehat j}.
 \label{eq:v8-horizontal-cofactor-vector}
\end{equation}
В физическом вакууме он равен
\begin{equation}
 c(A_0)=(0,0,0,0,0,0,-1,0,0,1,0)^{\mathsf T},
 \qquad A_0c(A_0)=0,
 \qquad c(A_0)^*c(A_0)=2.
 \label{eq:v8-horizontal-cofactor-kernel}
\end{equation}
Значит, determinant-line объект существует и канонически воспроизводит
одномерное ядро прямоугольной инцидентности.

На горизонтальной фазовой окружности
$A(z)=\operatorname{diag}(z^4I_3,z^3I_7)A_0$ каждый максимальный минор
получает determinant-характер строк:
\begin{equation}
 c(A(z))=z^{3\cdot4+7\cdot3}c(A_0)=z^{33}c(A_0).
 \label{eq:v8-horizontal-cofactor-phase-weight}
\end{equation}
Таким образом, кофакторная линия действительно видит плоскую фазу.

Однако фазочувствительность ещё не даёт допустимого скаляра. Суммарные
гиперзаряды determinant-характеров двух концов совпадают:
\begin{equation}
 Y_{\det E_s}=6\frac16+4\left(-\frac12\right)-1=-2,
 \qquad
 Y_{\det E_t}=3\frac23+3\left(-\frac13\right)-2+2\left(-\frac12\right)=-2.
 \label{eq:v8-horizontal-determinant-hypercharges}
\end{equation}
Их относительный характер нейтрален, поэтому носитель максимальных миноров
эквивалентен исходному 11-мерному модулю источника. Но этот модуль не
содержит тривиального представления:
\begin{equation}
 \dim\operatorname{Hom}_G(E_s,\mathbb C)=0.
 \label{eq:v8-horizontal-no-invariant-functional}
\end{equation}
Следовательно, калибровочно-инвариантной линейной свёртки кофакторного
вектора в скаляр нет.

Вещественное спаривание не помогает. Оно даёт только модуль
\begin{equation}
 c(A(z))^*c(A(z))=2,
 \label{eq:v8-horizontal-real-pair-modulus}
\end{equation}
не зависящий от $z$. Свёртка с фиксированным вакуумным вектором ядра была
бы фазочувствительной, но не канонической: уже два различных ковектора
$-e_7^*$ и $e_{10}^*$ дают одинаковую нормировку на $c(A_0)$. Такой выбор
повторяет прежнюю ошибку фиксированного фона.

\begin{theorem}[Фазочувствительная линия без скалярной тривиализации]
Кофакторная determinant-line несёт точный фазовый вес $33$ и порождает
ядро $A_0$, но её 11-мерный носитель не имеет калибровочно-инвариантного
линейного функционала. Полная Real-пара сохраняет только модуль.
\end{theorem}

\begin{nogocriterion}[Determinant-line no-go текущего носителя]
Нельзя поднимать горизонтальную моду членом
$\operatorname{Re}\langle\ell,c(A)\rangle$, пока ковектор $\ell$ не
выведен как эквивариантная секция. Выбор $\ell=c(A_0)^*$ вручную нарушает
фон-независимость, а Real-свёртка фазу сокращает. Ранние inflow- и
Pfaffian-гейты также не поставляют недостающую ориентацию.
\end{nogocriterion}

\begin{protocol}
\begin{enumerate}
 \item кофакторное ядро: \textbf{получено};
 \item фазовый вес determinant-line: \textbf{$33$};
 \item относительный determinant-гиперзаряд: \textbf{$0$};
 \item инвариантные линейные свёртки: \textbf{$0$};
 \item Real-пара сохраняет фазу: \textbf{нет};
 \item канонический фазочувствительный скаляр: \textbf{не получен};
 \item следующий гейт: \textbf{замыкание фазы тяжёлыми стрелочными циклами}.
\end{enumerate}
\end{protocol}

Машинный сертификат сохранён в
\path{s2t_v8_horizontal_phase_determinant_line_admission_gate_results.json}.